% !TEX root = ../main-es.tex
% Las cajas de ayuda no son contenido: texcount las ignora (ver `make wordcount`).
%TC:macro \limite [ignore,ignore]
%TC:macro \nota [ignore]
%TC:macro \TODO [ignore]
%TC:macro \parencite [ignore]
%TC:macro \textcite [ignore]
\section{Posibles aplicaciones}
\limite{0}{Posibles aplicaciones: plantee posibles aplicaciones prácticas del tema
trabajado, señalando su utilidad en contextos reales.}

\begin{description}[leftmargin=1.4em]

\item[Diseño de motores de ejecución agéntica.] Es la aplicación directa: en OASys, la
frontera medida en $(\lambda, c)$ decide si implementar la aplicación en caliente de
ediciones, con qué presupuesto reemplazar el máximo de iteraciones fijado a mano, y
cuándo la gobernanza (presupuesto $B$ más líneas inmutables) supera a la supervisión
humana por aprobación. La misma lógica aplica a orquestadores de agentes LLM
genéricos, que hoy fijan por intuición cuándo replanificar.

\item[Presupuestos de deliberación con LLM.] Modelar al modelo como un
oráculo de competencia $p$ y costo $c$ convierte una propiedad del modelo en un
parámetro del experimento: las curvas resultado dicen, para cualquier calidad y precio
del modelo, si la auto-edición de planes vale la pena. Un despliegue real calibra $p$
con unas pocas corridas y lee su decisión en la frontera ya construida.

\item[Replanificación en robótica y control.] La pregunta de cuándo replanificar bajo
dinamismo y observabilidad parcial es central en control predictivo de POMDPs
\parencite{ma2018performance} y en la simplificación de lazo abierto con garantías
formales \parencite{kong2026open}; los resultados de este proyecto ofrecen una versión
medible y reproducible de ese compromiso, útil como banco de pruebas docente.

\item[Gobernanza de agentes que se modifican a sí mismos.] Un agente que reescribe su
propio programa con líneas protegidas y presupuesto de ediciones es un caso mínimamente
viable del problema general de alinear un agente autor modificante: qué parte de sí
mismo no debería poder tocar, y cuántas veces conviene dejarlo intentarlo.

\end{description}
